\begin{proof} 1. \\
  \begin{itemize}
      \item[]\hspace{-0.5cm} ($(1) \rightarrow (2)$)\\
        Prop7.1.5.1 より明らか.
      \item[]\hspace{-0.5cm} ($(2) \rightarrow (1)$)\\
        対偶を示す.\\
        仮定とProp7.1.5.2より,
        $\BijMap {\seigen {C_A}{A}}{A}{\nat}$.\\
        Def7.1.4.2より$A$は可算無限集合である.\\
        よって対偶は真.
  \end{itemize}
\end{proof}


\begin{proof} 2. \\
  まず補題として以下の証明を行う.
  \begin{equation*}
   A \neq \emptyset \rightarrow \forall m,n \in A(C_A(m) \leq C_A(n) \leftrightarrow m \leq n )
  \end{equation*}
  $A \neq \emptyset$を仮定し,
  任意に$m, n \in A$をとる.\\
  \begin{itemize}
      \item[]\hspace{-0.5cm} ($m \leq n \rightarrow C_A(m) \leq C_A(n)$)\\
        前件を仮定する.\\
        $\varphi (n) :\Leftrightarrow C_A(m) \leq C_A(n)$とし,
        $n$に関する帰納法で示す.\\
        \begin{itemize}
          \item[]\hspace{-0.5cm} (base step)\\
            $C_A(n) = C_A(m)$より,
            $C_A(m) \leq C_A(n)$.\\
          \item[]\hspace{-0.5cm} (induction step)\\
            $C_A(n + 1) = C_A(n) \lor C_A(n) + 1$.
            よって,
            $C_A(n) \leq C_A(n+1)$\\
            これと$\varphi (n)$を合わせると,
            $C_A(m) \leq C_A(n+1)$
        \end{itemize}
      \item[]\hspace{-0.5cm} ($C_A(m) \leq C_A(n)  \rightarrow m \leq n$)\\
        $C_A(m) \leq C_A(n)$とする.\\
        $m \leq n \rightarrow C_A(m) \leq C_A(n)$と,
        $\seigen {C_A}{A}$は単射より以下は真.
          \begin{equation}
            \forall a,b \in A(a < b \rightarrow C_A(a) < C_A(b))
          \end{equation}
        \begin{sub-block}{{\boldmath $\because$}\,}
          \begin{itemize}
            \item[] \hspace{-0.5cm} \\
              任意に$a,b \in A$をとり$a < b$とする.\\
              $\seigen{C_A}{A}$は単射より,
              $\forall x, y \in A (x \neq y \rightarrow C_A(x) \neq C_A(y))$.\\
              $a,b \in A \land a \neq b$より$C_A(a) \neq C_A(b)$\\
              $m \leq n \rightarrow C_A(m) \leq C_A(n)$より,
              これに$a,b$を代入すると$C_A(a) \leq C_A(b)$.\\
              いま,
              $C_A(a) \neq C_A(b)$なので$C_A(a) < C_A(b)$.
          \end{itemize}
        \end{sub-block}
        命題2.9.1より,
        $\nat$は大小関係に関して全順序集合だから,
        以下は真.
        \begin{equation}
          \forall a, b \in \nat (a \leq b \lor b \leq a )
        \end{equation}
        これと(7.1)の対偶を合わせると,
        以下は真.
        \begin{equation}
            \forall a,b \in A(C_A(b) \leq C_A(a) \rightarrow b \leq a)
        \end{equation}
        \begin{sub-block}{{\boldmath $\because$}\,}
          \begin{itemize}
            \item[] \hspace{-0.5cm} \\
              (7.1)の対偶をとると,
              $\forall a, b, \in A(\lnot (C_A(a) < C_A(b)) \rightarrow \lnot (a < b))$.\\
              これに(7.2)を適用すると,
              $\forall a,b \in A(C_A(b) \leq C_A(a) \rightarrow b \leq a).$
          \end{itemize}
        \end{sub-block}
        $m,n \in A$,
        $C_A(m) \leq C_A(n)$と(7.3)により,
        $m \leq n$.
  \end{itemize}
  次に2を示す.\\
  $A \neq \emptyset$を仮定する.\\
  $l := \seigen{C_A}{A}^{-1}(0)$とおくと$l \in A \land C_A(l) = 0$.\\
  $0$は$\nat$の最小値なので
  $\forall x \in \nat (0 \leq x)$.
  よって$\forall n \in A (C_A(l) \leq C_A(n))$は真.\\
  $l,n \in A$,
  $C_A(l) \leq C_A(n)$と補題より,
  $l \leq n$.\\
  そんな$l$の存在より命題は真.
\end{proof}
